% =============================================================================
% 5‑Body.tex — State of the Art, General Description, Implementation, Results
% =============================================================================

% ╔═══════════════════════════════════════════════════════════════════════════╗
% ║  Chapter 2 — State of the Art                                           ║
% ╚═══════════════════════════════════════════════════════════════════════════╝
\chapter{State of the Art}
\label{ch:state-of-art}

\section{PID Control}
\label{sec:sota-pid}

The Proportional-Integral-Derivative (PID) controller is the most widely deployed feedback controller in industrial and educational robotics \cite{astrom2006}. Its output is a weighted sum of three terms: the proportional term reacts to the present error, the integral term eliminates steady-state offset by accumulating past errors, and the derivative term dampens oscillations by anticipating future error from the current rate of change \cite{ogata2010}. The standard formulation is:

\begin{equation}
    u(t) = K_P \cdot e(t) + K_I \int_0^t e(\tau)\,d\tau + K_D \frac{de(t)}{dt}
    \label{eq:pid}
\end{equation}

where $e(t)$ is the error (setpoint minus measured value) and $K_P$, $K_I$, $K_D$ are the tuning gains. The discrete-time implementation used in this project replaces the integral with a running sum and the derivative with a first-order backward difference, both gated by the time step $\Delta t$ to maintain physical consistency.

A known limitation of PID controllers is that the gains are fixed: the same $K_P$ applies whether the error is small or large. This produces a linear response characteristic that can cause oscillations near the setpoint (if gains are tuned for large errors) or sluggish correction far from the setpoint (if gains are tuned for stability).

\section{Fuzzy Logic and Fuzzy Control}
\label{sec:sota-fuzzy}

Fuzzy logic, introduced by Zadeh in 1965 \cite{zadeh1965}, generalizes Boolean logic by allowing truth values between 0 and 1. A crisp input value is mapped to one or more fuzzy sets through \textit{membership functions}, producing a degree of membership $\mu \in [0,1]$ that quantifies how well the input matches each linguistic term (e.g.\ ``Small'', ``Medium'', ``Large'').

A Mamdani-type fuzzy controller \cite{mamdani1975} operates in three stages:

\begin{enumerate}
    \item \textbf{Fuzzification:} the crisp input is converted into
     membership degrees using triangular or trapezoidal membership functions.
    \item \textbf{Inference:} a set of IF-THEN rules maps the fuzzy input to a fuzzy output. The AND operator is implemented as the minimum of the input memberships, and multiple rules activating the same output label are aggregated using the MAX operator.
    \item \textbf{Defuzzification:} the aggregated fuzzy output is converted back to a crisp value, typically using the centre-of-gravity (weighted average) method.
\end{enumerate}

Fuzzy PI controllers, which use fuzzy logic to adapt the gains of a conventional PI controller, have been applied to mobile robot navigation by several authors \cite{lee1990, passino1998}. The advantage over a fixed-gain PID is the ability to encode expert knowledge (``if the error is small, use gentle corrections'') without requiring a precise mathematical plant model.

\section{Depth-First Search for Maze Solving}
\label{sec:sota-dfs}

Depth-First Search (DFS) is a graph traversal algorithm that explores as far as possible along each branch before backtracking \cite{cormen2009}. Applied to a grid-based maze, DFS treats each cell as a node and each open passage as an edge. The algorithm maintains a stack of visited cells: when no unvisited neighbour is available, it pops the stack to return to the most recent junction with unexplored branches. DFS guarantees complete exploration of any connected graph and requires only $O(V)$ memory, where $V$ is the number of cells — well within the constraints of the EV3 MicroPython environment.


% ╔═══════════════════════════════════════════════════════════════════════════╗
% ║  Chapter 3 — General Description                                        ║
% ╚═══════════════════════════════════════════════════════════════════════════╝
\chapter{General Description}
\label{ch:general-description}

\section{System Overview}
\label{sec:system-overview}

The system consists of three layers: the \textit{hardware layer} (sensors, motors, EV3 brick), the \textit{control layer} (algorithms and task controllers), and the \textit{telemetry layer} (WiFi streaming to PC). Figure~\ref{fig:system-diagram} presents the high-level architecture.

\begin{figure}[htbp]
    \centering
    % ── Replace with your actual diagram ──
    \fbox{\parbox{0.85\textwidth}{\centering\vspace{4cm}
    \textit{[Insert system architecture diagram here]}\\[6pt]
    \small Sensors $\rightarrow$ Controller $\rightarrow$ Algorithm $\rightarrow$ Motors\\
    \hspace{4cm}$\downarrow$\\
    \hspace{4cm}WiFi Streamer $\rightarrow$ PC
    \vspace{1cm}}}
    \caption{High-level system architecture showing the data flow between the hardware layer, the control layer (with interchangeable algorithms), and the telemetry layer.}
    \label{fig:system-diagram}
\end{figure}

\section{Hardware Configuration}
\label{sec:hardware}

The robot is built from the LEGO Mindstorms EV3 Education Set. The drive train uses two large servomotors on ports A and C, connected through a differential gear to wheels with a diameter of 55.5\,mm and an axle track of 104\,mm. The Pybricks \texttt{DriveBase} class uses these dimensions to convert speed (mm/s) and turn rate (deg/s) commands into individual motor speeds.

For the lane keeping task, two colour sensors are mounted on ports S2 and S3, positioned symmetrically on the front of the robot to straddle the white lane between the two black boundary lines.

For the maze solving task, three ultrasonic sensors (forward on S4, left on S1, right on S2), a gyroscope (S3), and a sensor motor (B) that can rotate the forward sensor are used to perceive the maze and maintain accurate heading.

\section{Software Architecture}
\label{sec:software-arch}

The software is organized around the following abstractions:

\begin{itemize}
    \item \textbf{Configuration:} a single function \texttt{initializeConfiguration(task, algorithm)} generates a dictionary of all parameters for the selected task-algorithm pair.
    \item \textbf{Algorithm:} an abstract base class with a \texttt{compute(setpoint, measured, dt)} interface. Three concrete implementations (PID, Fuzzy PI, Fuzzy Rule-Based) are interchangeable at runtime.
    \item \textbf{Controller:} an abstract base class with a \texttt{run()} method. Two concrete controllers (LaneKeepingController, MazeSolverController) use the Algorithm interface to regulate the robot.
    \item \textbf{WifiStreamer:} a UDP socket wrapper that sends telemetry packets (step, disturbance, output) to a PC for real-time monitoring.
\end{itemize}

This architecture follows the SOLID principles: the Algorithm hierarchy satisfies the Open/Closed and Liskov Substitution principles (new algorithms can be added without modifying controllers), the single-responsibility principle separates sensing, computation and actuation, and the dependency inversion principle ensures that controllers depend on the abstract Algorithm interface rather than on concrete implementations.


% ╔═══════════════════════════════════════════════════════════════════════════╗
% ║  Chapter 4 — Implementation                                             ║
% ╚═══════════════════════════════════════════════════════════════════════════╝
\chapter{Implementation}
\label{ch:implementation}

% ─────────────────────────────────────────────────────────────────────────────
% 4.1  LANE KEEPING
% ─────────────────────────────────────────────────────────────────────────────
\section{The Lane Keeping Module}
\label{sec:lane-keeping}

\subsection{Task Description}
\label{sec:lk-task}

The lane keeping task requires the robot to follow a white surface bounded by two black lines. Two colour sensors mounted symmetrically on the front of the robot continuously measure the reflectance of the surface beneath them. Black surfaces return values close to 4, while white surfaces return values close to 14. The \textit{disturbance} is defined as the difference between the left and right sensor readings: a value of zero indicates perfect centering, while a positive or negative value indicates a drift to the right or left, respectively.

The task configuration defines the following thresholds:

\begin{table}[htbp]
\centering
\caption{Lane keeping configuration parameters.}
\label{tab:lk-config}
\begin{tabular}{l c p{7cm}}
\hline
\textbf{Parameter} & \textbf{Value} & \textbf{Description} \\
\hline
BlackThreshold      & 6       & Sensor value indicating the robot is on the black line. \\
CrossingThreshold   & 7       & Sensor value indicating the robot is very close to the black line. \\
MinDifference       & 3       & Tolerance zone for going straight (ignoring shadows). \\
BaseSpeed           & 90 mm/s & Forward velocity when centered. \\
TurningSpeed        & 70 mm/s & Reduced velocity during algorithm-regulated turns. \\
EmergencySpeed      & 65 mm/s & Restricted velocity during emergency corrections. \\
EmergencyTurnSpeed  & 100 deg/s & Extreme rotational speed for lane-exit recovery. \\
StopConfirmCount    & 1       & Consecutive black readings to confirm stop line. \\
\hline
\end{tabular}
\end{table}

\subsection{The LaneKeepingController}
\label{sec:lk-controller}

The \texttt{LaneKeepingController} class implements the \texttt{Controller} interface and orchestrates the four phases of the control cycle: sense, decide, actuate and wait.

The \texttt{decideAction()} method implements a three-tier safety-layered architecture:

\begin{enumerate}
    \item \textbf{Emergency override:} if one sensor reads below the crossing threshold while the other reads a normal value, the robot is about to leave the lane. The method immediately returns the emergency speed and maximum turn rate to aggressively pivot the robot back into the lane. This is a hard-coded safety net that overrides any algorithm output.
    \item \textbf{Dead-band tolerance:} if the absolute difference between the sensors is below the minimum difference threshold and the algorithm is PID or Fuzzy PI, the robot continues straight at base speed. This prevents the controller from reacting to sensor noise and shadows on the white surface. The dead-band is intentionally bypassed for the Fuzzy Rule-Based algorithm, which can handle small disturbances through its ZE (zero) fuzzy set.
    \item \textbf{Algorithm regulation:} if neither emergency nor dead-band conditions are met, the disturbance is passed to the selected algorithm's \texttt{compute()} method. The algorithm returns a turn rate, and the robot drives at the reduced turning speed with this computed correction.
\end{enumerate}

The \texttt{run()} method executes the continuous control loop. At each iteration it reads both sensors (\textit{sense}), calls \texttt{decideAction()} to determine speed and turn rate (\textit{decide}), sends the commands to the motors via \texttt{robot.drive()} (\textit{actuate}), and waits 20\,ms to maintain consistent loop frequency (\textit{wait}). A separate stop-line detection mechanism counts consecutive cycles where both sensors read below the black threshold; when this count reaches the confirmation threshold, the robot stops.


\subsection{PID Algorithm}
\label{sec:pid}

The \texttt{PIDAlgorithm} class implements the discrete PID controller from Equation~\ref{eq:pid}. The gains for the lane keeping task are $K_P = 0.8$, $K_I = 0.05$, $K_D = 0.003$. The output (turn rate) is clamped to $[-100, +100]$\,deg/s and the integral is clamped to $[-1000, +1000]$ to prevent windup.

The proportional term $K_P \cdot e$ provides immediate steering proportional to the current disturbance. The integral term $K_I \cdot \sum e \cdot \Delta t$ accumulates the disturbance over time to eliminate steady-state offset caused by mechanical drift (e.g.\ a slightly weaker motor). The derivative term $K_D \cdot \Delta e / \Delta t$ dampens oscillations by reducing the steering as the robot approaches the centre, preventing overshoot into the opposite boundary.


\subsection{Fuzzy PI Algorithm}
\label{sec:fuzzy-pi}

The \texttt{FuzzyPIAlgorithm} class uses fuzzy logic to dynamically adapt the proportional and integral gains of a PI controller. Unlike the fixed-gain PID, the Fuzzy PI adjusts its aggressiveness based on how far the robot has drifted from the centre: gentle corrections near the centre and aggressive corrections far from it.

\subsubsection{Fuzzification}

The \texttt{fuzzifyDisturbance()} method computes the degree of matching between the crisp absolute disturbance $|d|$ and three fuzzy linguistic terms: Small, Medium and Large. The membership functions are defined by three boundary values: $S = 3$, $M = 6$, $L = 10$.

\begin{itemize}
    \item \textbf{Small (trapezoidal):} $\mu_S = 1.0$ when $|d| \le S$; decreases linearly to 0 as $|d|$ increases from $S$ to $M$; zero when $|d| > M$.
    \item \textbf{Medium (triangular):} zero when $|d| \le S$ or $|d| \ge L$; rises from 0 to 1 between $S$ and $M$; falls from 1 to 0 between $M$ and $L$; peak at $|d| = M$.
    \item \textbf{Large (trapezoidal):} zero when $|d| \le M$; rises from 0 to 1 between $M$ and $L$; full membership when $|d| \ge L$.
\end{itemize}

Each formula on the transition slopes is derived from similar right triangles. For the descending slope of Small between $S$ and $M$, the full slope forms a right triangle with base $(M - S)$ and height 1. An input $|d|$ in this interval creates a smaller similar triangle with base $(M - |d|)$ and height $\mu_S$. By the ratio of corresponding sides:

\begin{equation}
    \mu_S(|d|) = \frac{M - |d|}{M - S}
    \label{eq:mu-small}
\end{equation}

The same geometric reasoning yields $\mu_M = (|d| - S)/(M - S)$ for the ascending slope, $\mu_M = (L - |d|)/(L - M)$ for the descending slope, and $\mu_L = (|d| - M)/(L - M)$ for the ascending slope of Large.

\subsubsection{Fuzzy Inference Rules}

The inference stage applies three IF-THEN rules that map each fuzzy set to gain multipliers:

\begin{table}[htbp]
\centering
\caption{Fuzzy PI inference rules.}
\label{tab:fuzzy-pi-rules}
\begin{tabular}{c l c l}
\hline
\textbf{Rule} & \textbf{IF disturbance is} & \textbf{THEN} & \textbf{Interpretation} \\
\hline
R1 & Small  & $K_P \times 0.8$,\; $K_I \times 0.7$ & Reduce gains near centre \\
R2 & Medium & $K_P \times 1.0$,\; $K_I \times 1.0$ & Nominal gains \\
R3 & Large  & $K_P \times 1.2$,\; $K_I \times 1.3$ & Boost gains far from centre \\
\hline
\end{tabular}
\end{table}

\subsubsection{Defuzzification and Output}

The adaptive multipliers are computed as weighted averages using the membership degrees as weights:

\begin{equation}
    K_{P,\text{mult}} = \frac{\mu_S \cdot 0.8 + \mu_M \cdot 1.0 + \mu_L \cdot 1.2}{\mu_S + \mu_M + \mu_L}
    \label{eq:kp-mult}
\end{equation}

The final adaptive gains are $K_{P,\text{adaptive}} = K_{P,\text{base}} \cdot K_{P,\text{mult}}$ and $K_{I,\text{adaptive}} = K_{I,\text{base}} \cdot K_{I,\text{mult}}$. The output is then computed as a standard PI controller with these adaptive gains.


\subsection{Fuzzy Rule-Based Algorithm}
\label{sec:fuzzy-rb}

The \texttt{FuzzyRuleBased} class eliminates the PI structure entirely and maps two fuzzy inputs — the disturbance and its rate of change (trend) — directly to a turn rate through a 15-rule Mamdani inference system.

\subsubsection{Fuzzification of Disturbance}

Unlike the Fuzzy PI which uses three sets on the absolute value, the Fuzzy Rule-Based uses five signed fuzzy sets: NL (Negative Large), NS (Negative Small), ZE (Zero), PS (Positive Small), PL (Positive Large). The membership functions are defined by $S = 3$ and $M = 6$:

\begin{itemize}
    \item \textbf{NL} (left shoulder): full at $d \le -M$, falls to 0 at $d = -S$.
    \item \textbf{NS} (triangle): peak at $d = -S$, base from $-M$ to $0$.
    \item \textbf{ZE} (triangle): peak at $d = 0$, base from $-S$ to $+S$.
    \item \textbf{PS} (triangle): peak at $d = +S$, base from $0$ to $+M$.
    \item \textbf{PL} (right shoulder): rises from 0 at $d = +S$, full at $d \ge +M$.
\end{itemize}

\subsubsection{Fuzzification of Trend}

The trend $\Delta d = d_{\text{current}} - d_{\text{previous}}$ is fuzzified into three sets with threshold $D = 3$:

\begin{itemize}
    \item \textbf{N} (Negative): full at $\Delta d \le -D$, falls to 0 at $\Delta d = 0$. The disturbance is recovering.
    \item \textbf{Z} (Zero): triangle, peak at $\Delta d = 0$, base from $-D$ to $+D$. The disturbance is stable.
    \item \textbf{P} (Positive): rises from 0 at $\Delta d = 0$, full at $\Delta d \ge +D$. The disturbance is worsening.
\end{itemize}

\subsubsection{Rule Table}

The 15 rules are organized in a $5 \times 3$ matrix. The output singletons are: NL$=-50$, NM$=-30$, NS$=-10$, ZE$=0$, PS$=+10$, PM$=+30$, PL$=+50$ deg/s.

\begin{table}[htbp]
\centering
\caption{Fuzzy Rule-Based: rule matrix for lane keeping.}
\label{tab:rule-matrix}
\begin{tabular}{l|c|c|c}
\hline
\textbf{Disturbance $\backslash$ Trend} & \textbf{N (recovering)} & \textbf{Z (stable)} & \textbf{P (worsening)} \\
\hline
NL (far left)      & PM (+30) & PL (+50) & PL (+50) \\
NS (slight left)   & PS (+10) & PM (+30) & PM (+30) \\
ZE (centered)      & ZE (0)   & ZE (0)   & ZE (0) \\
PS (slight right)  & NS ($-$10)  & NM ($-$30) & NM ($-$30) \\
PL (far right)     & NM ($-$30)  & NL ($-$50) & NL ($-$50) \\
\hline
\end{tabular}
\end{table}

The rule table has an antisymmetric structure: the output for (NL, Z) is PL $(+50)$ while the output for (PL, Z) is NL $(-50)$ — equal magnitude, opposite sign. The trend column modulates correction strength: when the robot is recovering (N), the correction is reduced to avoid overshoot; when worsening (P) or stable (Z), the full correction is applied.

\subsubsection{Inference, Aggregation and Defuzzification}

At each control cycle, up to two disturbance sets and two trend sets can have non-zero membership, activating 1, 2, or 4 rules simultaneously. The firing strength of each rule is computed as $\text{strength} = \min(\mu_{\text{dist}}, \mu_{\text{trend}})$. When multiple rules target the same output singleton, the MAX operator selects the highest firing strength. The crisp output is computed as a weighted average:

\begin{equation}
    \text{output} = \frac{\sum_i \text{agg}[i] \cdot \text{singleton}[i]}{\sum_i \text{agg}[i]}
    \label{eq:defuzz}
\end{equation}


% ─────────────────────────────────────────────────────────────────────────────
% 4.2  MAZE SOLVER
% ─────────────────────────────────────────────────────────────────────────────
\section{The Maze Solver Module}
\label{sec:maze-solver}

\subsection{Task Description}
\label{sec:ms-task}

The maze solving task requires the robot to navigate through an unknown grid-based maze composed of rectangular cells (260\,mm $\times$ 280\,mm) separated by walls. The robot must explore every reachable cell using a Depth-First Search strategy, then return to the starting position. Between junctions the same control algorithms regulate the forward approach speed.

\begin{table}[htbp]
\centering
\caption{Maze solver configuration parameters.}
\label{tab:ms-config}
\begin{tabular}{l c p{6.5cm}}
\hline
\textbf{Parameter} & \textbf{Value} & \textbf{Description} \\
\hline
CellDistanceNS        & 260 mm   & North-South cell dimension. \\
CellDistanceEW        & 280 mm   & East-West cell dimension. \\
OpenPassageThreshold   & 150 mm   & Sensor $>$ this $\Rightarrow$ passage is open. \\
ObstacleThreshold      & 80 mm    & Forward $<$ this $\Rightarrow$ hard stop. \\
TargetForwardDistance  & 150 mm   & Algorithm setpoint for speed regulation. \\
BaseSpeed / MinSpeed   & 150 / 30 mm/s & Speed bounds during cell driving. \\
CurveSpeed / CurveTurnRate & 80 mm/s / 15 deg/s & Parameters for gyro-guided 90\textdegree{} turns. \\
BackupDistance         & 60 mm    & Reverse distance before turning. \\
HeadingThreshold       & 3\textdegree{}    & Dead zone for gyro correction. \\
CenteringThreshold / Gain & 20 mm / 0.1 & Side-sensor corridor centering. \\
\hline
\end{tabular}
\end{table}


\subsection{The MazeMap Class}
\label{sec:mazemap}

The \texttt{MazeMap} class maintains a coordinate-based map of the explored maze and implements the DFS navigation logic. It uses a compass-based coordinate system where the four cardinal directions are encoded as integers: 0 = North, 1 = East, 2 = South, 3 = West.

\begin{table}[htbp]
\centering
\caption{Direction encoding and coordinate changes.}
\label{tab:directions}
\begin{tabular}{l c c c c}
\hline
\textbf{Direction} & \textbf{Index} & \textbf{DX} & \textbf{DY} & \textbf{Gyro Target} \\
\hline
North & 0 & 0  & $+1$ & 0\textdegree{} \\
East  & 1 & $+1$ & 0  & 90\textdegree{} \\
South & 2 & 0  & $-1$ & 180\textdegree{} \\
West  & 3 & $-1$ & 0  & 270\textdegree{} \\
\hline
\end{tabular}
\end{table}

The DFS exploration follows the priority order Forward $\rightarrow$ Left $\rightarrow$ Right $\rightarrow$ Behind. The \texttt{planMove()} method iterates this priority, checking each candidate direction for two conditions: the direction must be physically open (detected by the ultrasonic sensors) and the neighbouring cell must not have been previously visited. The \texttt{planBacktrack()} method pops the path stack to return to the most recent junction with unexplored branches.

The navigation uses a two-phase commit pattern: \texttt{applyTurn()} updates the heading before the physical turn, then \texttt{commitMove()} pushes the current position onto the path stack, advances the coordinates, and marks the new cell as visited. When backtracking, \texttt{commitBacktrack()} pops the previous position from the stack.


\subsection{The MazeSolverController}
\label{sec:ms-controller}

The \texttt{MazeSolverController} orchestrates the physical execution of the DFS algorithm through five support routines:

\textbf{Sensor Reading.} The \texttt{scanForward()} method takes three consecutive readings, sorts them, and returns the median to reject ultrasonic spike noise. The \texttt{scanJunction()} method reads all three ultrasonic sensors and returns the set of open absolute directions.

\textbf{Heading Alignment.} After every turn, the \texttt{alignHeading()} method corrects the robot's physical orientation using the gyroscope. The error between the expected gyro target and the actual reading is normalized to $[-180^\circ, +180^\circ]$ and the robot executes a corrective point turn if the error exceeds 3\textdegree{}.

\textbf{Corridor Centering.} The \texttt{centerInCorridor()} method reads both side sensors and, when both walls are visible, computes a corrective turn proportional to the lateral offset ($0.1$\,deg/mm, capped at $\pm$15\textdegree{}).

\textbf{Turn Execution.} The \texttt{executeTurn()} method performs gyro-guided curve turns for 90\textdegree{} rotations (smooth arc at 80\,mm/s, 15\,deg/s) and encoder-based point turns for 180\textdegree{} U-turns.

\textbf{Cell-to-Cell Driving.} The \texttt{\_driveCell()} method drives one cell using the selected algorithm to regulate forward speed. The distance is tracked via trapezoidal integration: $s \mathrel{+}= \frac{1}{2}(v_{\text{prev}} + v_{\text{curr}}) \cdot \Delta t$. The loop exits when the integrated distance reaches the cell dimension or when the forward sensor detects an unexpected wall.


\subsection{Main Execution Loop}

The \texttt{run()} method repeats a four-phase cycle for each cell:

\begin{enumerate}
    \item \textbf{Sense:} scan the junction with \texttt{scanJunction()}.
    \item \textbf{Plan:} call \texttt{planMove()} or \texttt{planBacktrack()}.
    \item \textbf{Execute:} backup $\rightarrow$ turn $\rightarrow$ align heading $\rightarrow$ centre in corridor $\rightarrow$ drive cell $\rightarrow$ commit.
    \item \textbf{Repeat} until the path stack is empty (maze fully explored).
\end{enumerate}


% ─────────────────────────────────────────────────────────────────────────────
% 4.3  THE TELEMETRY INTERFACE
% ─────────────────────────────────────────────────────────────────────────────
\section{The Telemetry Interface}
\label{sec:ui}

\subsection{WifiStreamer}
\label{sec:wifi-streamer}

The \texttt{WifiStreamer} class provides real-time telemetry by streaming control data from the EV3 to a PC over a UDP socket. At initialization, it creates a \texttt{SOCK\_DGRAM} socket and sends a metadata packet containing the task-algorithm label. During the control loop, each cycle sends a comma-separated packet with the step number, the disturbance and the algorithm output.

UDP was chosen over TCP because it is connectionless and does not block the control loop if the PC is temporarily unreachable. The trade-off is that packets may be lost, but at the 50\,Hz control rate, occasional packet loss is acceptable for visualization purposes.

\subsection{PC Receiver and Visualization}
\label{sec:pc-receiver}

On the PC side, a Python script listens on the configured port (5005) and logs incoming packets to a CSV file. The CSV data is then visualized using Matplotlib to generate the comparative plots presented in Chapter~\ref{ch:results}. For the maze solver, the streamed $(x, y)$ cell coordinates are plotted on a grid to visualize the exploration path in real time.

The metadata packet sent at the beginning of each run identifies the task and algorithm, allowing the PC receiver to label the log files automatically (e.g.\ \texttt{LaneKeeping-PID.csv}, \texttt{MazeSolver-FuzzyRuleBased.csv}).


% ─────────────────────────────────────────────────────────────────────────────
% 4.4  OOP ARCHITECTURE
% ─────────────────────────────────────────────────────────────────────────────
\section{Object-Oriented Architecture}
\label{sec:oop}

\subsection{Class Hierarchy}
\label{sec:class-hierarchy}

The software is organized around two parallel hierarchies connected by composition:

\begin{figure}[htbp]
    \centering
    \fbox{\parbox{0.85\textwidth}{\centering\vspace{3cm}
    \textit{[Insert UML class diagram here]}\\[6pt]
    \small ControlStrategy $\leftarrow$ Algorithm $\leftarrow$ \{PID, FuzzyPI, FuzzyRuleBased, FuzzyStateMachine\}\\
    Controller $\leftarrow$ \{LaneKeepingController, MazeSolverController\}\\
    Controller \textit{has-a} Algorithm, WifiStreamer
    \vspace{1cm}}}
    \caption{UML class diagram showing the Algorithm and Controller hierarchies.}
    \label{fig:uml}
\end{figure}

The \textbf{Algorithm hierarchy} defines the control strategies. The abstract base class \texttt{Algorithm} declares the \texttt{compute(setpoint, measured, dt)} interface. Three concrete subclasses implement this interface:

\begin{itemize}
    \item \texttt{PIDAlgorithm}: fixed-gain PID controller.
    \item \texttt{FuzzyPIAlgorithm}: fuzzy-adaptive PI controller with 3 rules.
    \item \texttt{FuzzyRuleBased}: pure Mamdani controller with 15 rules and 2 inputs.
\end{itemize}

The \textbf{Controller hierarchy} defines the task-specific execution logic. The abstract base class \texttt{Controller} provides the stopwatch, step counter, and WiFi streamer. Two concrete subclasses implement the \texttt{run()} method:

\begin{itemize}
    \item \texttt{LaneKeepingController}: continuous control loop with emergency override.
    \item \texttt{MazeSolverController}: DFS-based exploration with cell-to-cell driving.
\end{itemize}

\subsection{SOLID Principles}
\label{sec:solid}

The architecture satisfies the five SOLID principles:

\textbf{Single Responsibility:} each class has one responsibility. \texttt{PIDAlgorithm} computes PID output; \texttt{MazeMap} manages graph traversal; \texttt{WifiStreamer} handles network communication.

\textbf{Open/Closed:} new algorithms can be added by creating a new subclass of \texttt{Algorithm} without modifying any existing code. The \texttt{FuzzyRuleBased} algorithm was added months after the initial PID implementation, requiring zero changes to the controllers.

\textbf{Liskov Substitution:} any \texttt{Algorithm} subclass can replace any other without breaking the controllers, because they all satisfy the same \texttt{compute()} contract.

\textbf{Interface Segregation:} the \texttt{Algorithm} interface exposes only \texttt{compute()} and \texttt{reset()}, rather than forcing all algorithms to implement fuzzification methods they don't need.

\textbf{Dependency Inversion:} the controllers depend on the abstract \texttt{Algorithm} interface. The concrete algorithm is selected at startup by the factory function \texttt{createController()}.

\subsection{Factory Pattern}
\label{sec:factory}

The \texttt{createController(task, algorithm\_name)} function serves as a factory that instantiates the correct algorithm and controller based on the configuration strings. It calls \texttt{initializeConfiguration()} to generate the parameter dictionary, \texttt{hardware\_setup()} to initialize the motors, and then creates the appropriate \texttt{Algorithm} and \texttt{Controller} objects. This centralizes the creation logic and ensures that the configuration, algorithm and controller are always consistent.


% ╔═══════════════════════════════════════════════════════════════════════════╗
% ║  Chapter 5 — Experimental Results                                        ║
% ╚═══════════════════════════════════════════════════════════════════════════╝
\chapter{Experimental Results}
\label{ch:results}

\section{Experimental Setup}
\label{sec:exp-setup}

All experiments were conducted on the same physical track (for lane keeping) and the same physical maze (for maze solving) to ensure fair comparison between the three algorithms. Each algorithm was run five times, and the telemetry data was averaged. The metrics used for evaluation are:

\begin{itemize}
    \item \textbf{Mean Absolute Disturbance (MAD):} the average of $|d|$ across all control steps, measuring how far the robot strays from the ideal state.
    \item \textbf{Standard Deviation of Disturbance ($\sigma_d$):} measures oscillation amplitude.
    \item \textbf{Crossing Count:} the number of times the emergency override was triggered.
    \item \textbf{Average Speed:} the mean forward velocity, reflecting how confidently the robot drives.
    \item \textbf{Completion Time:} total time to complete the track or maze.
\end{itemize}


\section{Lane Keeping Results}
\label{sec:lk-results}

\subsection{Disturbance Over Time}

Figure~\ref{fig:lk-disturbance} shows the disturbance $d$ plotted against control steps for the three algorithms on the same track.

\begin{figure}[htbp]
    \centering
    \fbox{\parbox{0.85\textwidth}{\centering\vspace{5cm}
    \textit{[Insert disturbance vs.\ step plot for PID, FuzzyPI, FuzzyRuleBased]}
    \vspace{1cm}}}
    \caption{Disturbance over time for the three algorithms in the lane keeping task.}
    \label{fig:lk-disturbance}
\end{figure}

\subsection{Performance Comparison}

\begin{table}[htbp]
\centering
\caption{Lane keeping performance comparison (average of 5 runs).}
\label{tab:lk-performance}
\begin{tabular}{l c c c c c}
\hline
\textbf{Algorithm} & \textbf{MAD} & $\boldsymbol{\sigma_d}$ & \textbf{Crossings} & \textbf{Avg Speed} & \textbf{Time} \\
\hline
PID             & --- & --- & --- & --- mm/s & --- s \\
Fuzzy PI        & --- & --- & --- & --- mm/s & --- s \\
Fuzzy Rule-Based & --- & --- & --- & --- mm/s & --- s \\
\hline
\end{tabular}
\end{table}

% ── Placeholder for analysis paragraphs ──
The results show that \textit{[fill with analysis after collecting data: which algorithm had the lowest MAD, fewest crossings, smoothest trajectory, and why the trend input in the Fuzzy Rule-Based reduces oscillations]}.


\subsection{Turn Rate Output}

Figure~\ref{fig:lk-output} compares the turn rate output of the three algorithms.

\begin{figure}[htbp]
    \centering
    \fbox{\parbox{0.85\textwidth}{\centering\vspace{5cm}
    \textit{[Insert turn rate output vs.\ step plot for all three algorithms]}
    \vspace{1cm}}}
    \caption{Turn rate output over time for the three algorithms.}
    \label{fig:lk-output}
\end{figure}

\subsection{Rule Activation Analysis}

For the Fuzzy Rule-Based algorithm, Figure~\ref{fig:rule-activation} shows the distribution of rule activations across the 15 rules during a typical run.

\begin{figure}[htbp]
    \centering
    \fbox{\parbox{0.85\textwidth}{\centering\vspace{5cm}
    \textit{[Insert rule activation histogram or heatmap]}
    \vspace{1cm}}}
    \caption{Rule activation frequency for the Fuzzy Rule-Based controller.}
    \label{fig:rule-activation}
\end{figure}


\section{Maze Solver Results}
\label{sec:ms-results}

\subsection{Exploration Map}

Figure~\ref{fig:maze-map} shows the cell-by-cell exploration path reconstructed from the streamed $(x, y)$ coordinates.

\begin{figure}[htbp]
    \centering
    \fbox{\parbox{0.85\textwidth}{\centering\vspace{5cm}
    \textit{[Insert maze exploration path visualization]}
    \vspace{1cm}}}
    \caption{Maze exploration path showing the DFS traversal order.}
    \label{fig:maze-map}
\end{figure}

\subsection{Speed Regulation}

Figure~\ref{fig:ms-speed} shows the forward speed during a cell-to-cell drive for each algorithm.

\begin{figure}[htbp]
    \centering
    \fbox{\parbox{0.85\textwidth}{\centering\vspace{5cm}
    \textit{[Insert speed vs.\ distance plot for one cell drive]}
    \vspace{1cm}}}
    \caption{Forward speed regulation during cell-to-cell driving.}
    \label{fig:ms-speed}
\end{figure}

\subsection{Maze Solver Performance}

\begin{table}[htbp]
\centering
\caption{Maze solver performance comparison.}
\label{tab:ms-performance}
\begin{tabular}{l c c c}
\hline
\textbf{Algorithm} & \textbf{Cells Explored} & \textbf{Total Time} & \textbf{Heading Error (avg)} \\
\hline
PID             & --- & --- s & ---\textdegree{} \\
Fuzzy PI        & --- & --- s & ---\textdegree{} \\
Fuzzy Rule-Based & --- & --- s & ---\textdegree{} \\
\hline
\end{tabular}
\end{table}
